\documentclass[10pt,a4paper]{amsart}
\usepackage[margin=19mm]{geometry}
\usepackage{amsmath,amssymb,mathtools}
\usepackage[scr=rsfs,cal=euler]{mathalfa}
\usepackage{xfrac}
\usepackage[unicode,hidelinks,bookmarks=false]{hyperref}
\newcommand{\ie}{i.e.\ }
\newcommand{\cf}{cf.\ }
\newcommand{\resp}{respectively\ }
\newcommand{\Subsection}{\subsection}
\providecommand{\nobreakdash}{\nobreak\hbox{-}}
\setlength{\emergencystretch}{2em}
\multlinegap0pt
\usepackage{amssymb}
\usepackage{tikz}
\usepackage{float}
\usepackage[shortlabels]{enumitem}
\usepackage{algorithm}\usepackage{algpseudocode}
\newtheorem{lemma}{Lemma}
\newcommand{\cA}{\mathcal{A}}
\newcommand{\cC}{\mathcal{C}}
\newcommand{\cT}{\mathcal{T}}
\newcommand{\cbra}[1]{\left\{#1\right\}}
\newcommand{\sC}{\mathsf{C}}
\newcommand{\sD}{\mathsf{D}}
\newcommand{\sQ}{\mathsf{Q}}
\newcommand{\sR}{\mathsf{R}}
\title{Complexity of learning matchings and half graphs via edge queries}
\author{Nikhil S. Mande, Swagato Sanyal, Viktor Zamaraev}
\date{Source: arXiv:2507.03151v2 (CC BY 4.0)}
\begin{document}
\maketitle

\noindent\emph{Source and licence.} Complexity of learning matchings and half graphs via edge queries. Version arXiv:2507.03151v2 (CC BY 4.0), distributed by arXiv under the Creative Commons Attribution 4.0 International licence, http://creativecommons.org/licenses/by/4.0/, as recorded in the arXiv licence metadata for this exact version and shown on the abstract page https://arxiv.org/abs/2507.03151v2. Full licence notice: \texttt{licenses/arxiv-2507.03151-CC-BY-4.0.txt}.\par\smallskip

\noindent\textit{Editorial scope.} Counted unit: the article's lemma lem:sorting-thr-equiv-colpermuted (source0.tex lines 525--530). The article's complete original proof of it is delivered with it (source0.tex lines 531--542, one \texttt{proof} environment of 12 lines). No mathematics is written, repaired or re-derived: the statement and the proof are the article's own lines, in the article's own notation, and no retained line is substituted or re-flowed.  No further source-local context is needed: every cross-reference of every retained line resolves inside the two delivered spans.  Nothing was omitted from any retained span, so the package carries no omission record.  No retained line contains a citation, so the wrapper carries no bibliography and no bibliography entry.  No retained line contains a figure environment, an image-inclusion command or drawing code, so no image is delivered and the package declares no asset dependency.  Editorial formatting only, in addition: an AMS \texttt{amsart} preamble that restates the article's own declarations and macros used by the retained lines, namely the environments \texttt{lemma} on separate counters, together with the standard packages those lines need.  The retained environments print in this document as Lemma 1 in order of appearance, whereas the article prints the same items with the numbering of its own full document.  The complete native source of the article as distributed in the arXiv e-print for this exact version is bundled unchanged as \texttt{sources/180999/source0.tex}.
\begin{lemma}\label{lem:sorting-thr-equiv-colpermuted}
    Let $n$ be a positive integer. Then for all $\sC \in \cbra{\sD, \sR, \sQ}$,
    $$
    \sC^t(\cT_n) = \sC(\cC_n).
    $$
\end{lemma}

\begin{proof}
    To prove the lemma, we establish a bijection between the lists in $\cT_n$ and matrices in $\cC_n$, and show that threshold queries to a list in $\cT_n$ are equivalent to edge queries to the corresponding matrix in $\cC_n$.

    The bijection is defined as follows. To a list $X \in \cT_n$ we associate the matrix $M_X \in \cC_n$, where the $j$-th column of $M_X$ is $0^{n - X[j]}1^{X[j]}$. That is, the top $n - X[j]$ entries are 0s, and the remaining entries are 1s. According to this bijection, a matrix $M \in \cC_n$ corresponds to the list $X_M \in \cT_n$, where for every $j \in [n]$, $X_M[j]$ is equal to the number of 1s in the $j$-th column of $M$.
    To show the equivalence between queries, we observe that for any $X \in \cT_n$ and $M \in \cC_n$, we have 
    $X[j] \geq n - i + 1$ if and only if $M_X(i, j) = 1$, and
    $M(i, j) = 1$ if and only if $X_M[j] \geq n - i + 1$.
    
    Thus, any algorithm $\cA$ for identifying a list in $\cT_n$ can be turned into an algorithm, with the exact same number of queries, for identifying a matrix in $\cC_n$ by replacing every query ``Is $X[j] \geq i$?'' to the hidden list $X$ with the query ``Is $M(n-i+1,j) = 1$?'' to the hidden matrix $M$.

    The argument to show that an algorithm for identifying a matrix in $\cC_n$ can be turned into an algorithm for identifying a list in $\cT_n$ is similar, and we omit it.
\end{proof}

\end{document}
